% VI32-DETAILED-PROBLEM-05
\begin{problem}[VI.32.P5: Nakayama and embedding dimension]\label{prob:vi32-05}
For a Noetherian local ring \((R,\mathfrak m,\kappa)\), prove that
\[
\dim_\kappa\mathfrak m/\mathfrak m^2
\]
is the minimal number of generators of \(\mathfrak m\).

\textbf{Provenance.} Canonical VI/32 synthesis.
\end{problem}

\ifdefined\FullProblemDossiers
\begin{solution}
Suppose first that
\[
\mathfrak m=(a_1,\ldots,a_r).
\]
Passing to the quotient by \(\mathfrak m^2\), the classes
\[
\bar a_1,\ldots,\bar a_r
\]
span the \(\kappa\)-vector space \(\mathfrak m/\mathfrak m^2\). Therefore
\[
\dim_\kappa\mathfrak m/\mathfrak m^2\le r.
\]
So the cotangent dimension is a lower bound for the number of generators.

Conversely, choose elements \(a_1,\ldots,a_e\in\mathfrak m\) whose classes form a basis of
\(\mathfrak m/\mathfrak m^2\), where
\[
e=\dim_\kappa\mathfrak m/\mathfrak m^2.
\]
Then
\[
\mathfrak m=(a_1,\ldots,a_e)+\mathfrak m^2.
\]
Let
\[
M=\mathfrak m/(a_1,\ldots,a_e).
\]
The displayed equality says
\[
M=\mathfrak mM.
\]
Because \(R\) is Noetherian, \(\mathfrak m\) and hence \(M\) are finitely generated. Nakayama's lemma therefore
gives \(M=0\). Thus
\[
\mathfrak m=(a_1,\ldots,a_e).
\]
No smaller generating set can exist, since its images would have to span an \(e\)-dimensional vector space.
Therefore
\[
\boxed{
e=\dim_\kappa\mathfrak m/\mathfrak m^2
=
\text{minimal number of generators of }\mathfrak m.
}
\]
\end{solution}
\fi
